\documentclass[a4paper,12pt]{article}
\usepackage{../packages/coursCollege}
\newcommand{\Chapitre}{Trigonométrie du triangle rectangle}
\renewcommand{\path}{../../}
\begin{document}

\section{Activités d'introduction}

\begin{activite}
	\tcblower
Sur la feuille distribuée en annexe sont tracés neuf triangles. Les mesures de leurs côtés sont indiqués en centimètre (ce sont les longueurs mesurées). 
\begin{tasks}
\task Calculer pour chaque triangle les rapports suivants par rapport à l'angle indiqué : 
\begin{gather*}
	\dfrac{\text{longueur du côté opposé}}{\text{longueur de l'hypoténuse}}\\\\
	\dfrac{\text{longueur du côté opposé}}{\text{longueur de l'hypoténuse}}\\\\
	\dfrac{\text{longueur du côté opposé}}{\text{longueur de l'hypoténuse}}
\end{gather*}
\task Que constatez-vous~? Comment peut-on expliquer~?
\task Soit un triangle $\triangle ABC$ rectangle en $C$. Calculer $\overline{BC}$ lorsque l'hypoténuse mesure $16$ cm et $\alpha=40^\circ$.
\end{tasks}
\end{activite}
\begin{activite}
\begin{center}
	\begin{tikzpicture}[scale=0.5]
		  % Définition des points
  \tkzDefPoint(0,0){A}
  \tkzDefPoint(4,0){B}
  \tkzDefPoint(4,3){C}
  \tkzDrawSegment(A,C)
  \tkzDrawSegment(A,B)
  \tkzDrawSegment(B,C)
  % Marquage d'un des angles opposés par le sommet
  \tkzMarkAngle[size=0.8cm, mark=none](B,A,C)
  \tkzLabelAngle[pos=1.2](B,A,C){$\alpha$}
  \tkzMarkRightAngle[size=0.5](C,B,A)
  \tkzLabelPoints(A,B)
  \tkzLabelPoints[above](C)
	\end{tikzpicture}
	\hspace{2cm}
	\begin{tikzpicture}[scale=0.5]
		  % Définition des points
  \tkzDefPoint(0,0){D}
  \tkzDefPoint(6,0){E}
  \tkzDefPoint(6,4.5){F}
  \tkzDrawSegment(D,F)
  \tkzDrawSegment(D,E)
  \tkzDrawSegment(E,F)
  % Marquage d'un des angles opposés par le sommet
  \tkzMarkAngle[size=0.8cm, mark=none](E,D,F)
  \tkzLabelAngle[pos=1.2](E,D,F){$\alpha$}
  \tkzMarkRightAngle[size=0.5](F,E,D)
  \tkzLabelPoints(D,E)
  \tkzLabelPoints[above](F)
	\end{tikzpicture}
\end{center}
{\bfseries Indice :} utiliser les termes côté adjacent à $\alpha$, côté opposé à $\alpha$ et hypoténuse.
	\tcblower 
	\begin{tasks}
		\task Expliquer pourquoi des triangles rectangles ayant un autre angle isométrique sont semblables. 
		\task Utiliser le théorème de Thalès pour compléter les rapports suivants~:

		\[\dfrac{\overline{BC}}{\overline{EF}}=\dfrac{\phantom{aaaa}}{\phantom{aaaa}}=\dfrac{\phantom{aaaa}}{\phantom{aaaa}}\]

		\task Compléter à présent les énoncés suivants~:

		Grâce à l'égalité $\dfrac{\phantom{aaaa}}{\phantom{aaaa}}=\dfrac{\phantom{aaaa}}{\phantom{aaaa}}$, on peut écrire \[\dfrac{\overline{BC}}{\overline{AC}}=\dfrac{\phantom{aaaa}}{\phantom{aaaa}}=\dfrac{\phantom{aaaaaaaaaaaaaa}}{\phantom{aaaaaaaaaaaaaaaaaaaaa}}\]

		Grâce à l'égalité $\dfrac{\phantom{aaaa}}{\phantom{aaaa}}=\dfrac{\phantom{aaaa}}{\phantom{aaaa}}$, on peut écrire \[\dfrac{\overline{BC}}{\overline{AB}}=\dfrac{\phantom{aaaa}}{\phantom{aaaa}}=\dfrac{\phantom{aaaaaaaaaaaaaa}}{\phantom{aaaaaaaaaaaaaaaaaaaaa}}\]

		Grâce à l'égalité $\dfrac{\phantom{aaaa}}{\phantom{aaaa}}=\dfrac{\phantom{aaaa}}{\phantom{aaaa}}$, on peut écrire \[\dfrac{\overline{AB}}{\overline{AC}}=\dfrac{\phantom{aaaa}}{\phantom{aaaa}}=\dfrac{\phantom{aaaaaaaaaaaaaa}}{\phantom{aaaaaaaaaaaaaaaaaaaaa}}\]
	\end{tasks}
\end{activite}
\newpage
\begin{landscape}
\pagestyle{empty}
\tikzset{
  myrotate/.style={rotate=#1,nodes={rotate=#1}}
}
\begin{tikzpicture}[remember picture, overlay,x=1cm, y=1cm]

%-----------------------------------------------------------------
% Helper macro to draw a triangle with sides (a,b,c) 
% placing AB=a horizontally, marking angle at A.
% Usage: \DrawTriangle{a}{b}{c}
% The code places A at (0,0) and B at (a,0), then computes C 
% using the law of cosines so that AC=b and BC=c.
% Also labels each side with its length and marks \angle A.

\newcommand{\DrawTriangle}[3]{%
  % #1=a, #2=b, #3=c
  \tkzDefPoint(0,0){A}
  \tkzDefPoint(#1,0){B}
  % Compute coordinates of C via law of cosines:
  %  x = (a^2 + b^2 - c^2) / (2a)
  %  y = sqrt(b^2 - x^2)
  \pgfmathsetmacro{\xC}{(\the\numexpr#1 * #1 + #2 * #2 - #3 * #3) / (2 * #1)} 
  \pgfmathsetmacro{\tempA}{(#2 * #2) - (\xC * \xC)} % inside of sqrt
  % guard against tiny floating remainders:
  \ifdim \tempA pt < 0.000001pt
    \pgfmathsetmacro{\tempA}{0}%
  \fi
  \pgfmathsetmacro{\yC}{sqrt(\tempA)}
  \tkzDefPoint(\xC,\yC){C}

  % Draw and label sides
  \tkzDrawSegments(A,B B,C C,A)
  \tkzLabelSegment[sloped](A,B){\pgfmathprintnumber{#1} cm}
  % For the other two sides, you can adjust label placement if needed:
  \tkzLabelSegment[sloped](A,C){\pgfmathprintnumber{#2} cm}
  \tkzLabelSegment[sloped](B,C){\pgfmathprintnumber{#3} cm}

  % Mark the angle at A
  % (You can adjust "size=..." or "mark=..." as you prefer)
  \tkzMarkAngle[size=0.8](B,A,C)
}

%-----------------------------------------------------------------
% Now place each triangle in its own "scope" with some shift 
% so they don’t overlap. We have 9 triangles total.

% The triangles you gave, in the order you gave them:
% 1) 7.84, 5.04, 6
% 2) 8.54, 7.4, 4.27
% 3) 12.41, 9.5, 7.97
% 4) 13.29, 9.4, 9.4
% 5) 6.7, 9.48, 6.7
% 6) 4.5, 3.9, 2.25
% 7) 5.66, 4, 4
% 8) 10.39, 9, 5.2
% 9) 5, 3.83, 3.21

% Adjust these shifts so everything fits on the page
% (You can also create multiple pages or bigger page sizes.)
% Example: 3 columns, 3 rows
\begin{scope}[shift={(22,-5)}, myrotate=90]
  \DrawTriangle{4.95}{3.83}{3.21}
\end{scope}

\begin{scope}[shift={(18,-17.5)}, myrotate=40]
  \DrawTriangle{7.84}{6}{5.04}
\end{scope}

\begin{scope}[shift={(17,-18)}, myrotate=110]
  \DrawTriangle{8.54}{7.4}{4.27}
\end{scope}

\begin{scope}[shift={(-1.5,-17.5)}, rotate=0]
  \DrawTriangle{12.41}{9.5}{7.97}
\end{scope}

\begin{scope}[shift={(14.5,-5.5)}, rotate=180]
  \DrawTriangle{13.29}{9.4}{9.4}
\end{scope}

\begin{scope}[shift={(15,-6)}, myrotate=320]
  \DrawTriangle{6.7}{9.48}{6.7}
\end{scope}

\begin{scope}[shift={(3,-13)}, myrotate=120]
  \DrawTriangle{4.5}{3.9}{2.25}
\end{scope}

\begin{scope}[shift={(15,0)}, myrotate=230]
  \DrawTriangle{5.66}{4}{4}
\end{scope}

\begin{scope}[shift={(-1,-4)}]
  \DrawTriangle{10.39}{9}{5.2}
\end{scope}
\end{tikzpicture}

 
\end{landscape}
\end{document}

